\section{绪论}
\subsection{研究背景}
操作系统是现代计算机体系结构的核心组成部分，承担着管理与配置计算资源与支持上层软件运行的关键职责。近年来，随着多核处理器规模的持续扩展以及专用加速器的广泛应用，计算体系结构正由传统的对称多处理（SMP）模式逐步转向异构计算架构。Jia 等人（2025）指出，随着 CPU、GPU 计算能力的提升，以及专为 AI 计算设计的 NPU 的出现与普及，异构计算已成为当前主流趋势\cite{jia2024agentcentricoperating}。

与此同时，计算负载形态也发生显著变化。云计算与数据中心环境下的任务具有高度动态性与多样性，微服务架构与容器化部署进一步增加了系统调度的复杂性。AI 训练与推理任务的规模化运行，使得资源需求呈现出高维、多阶段与强耦合特征。

在硬件结构与负载形态双重变化的驱动下，操作系统内部的资源调度问题逐渐从相对稳定、低维的分配模型转向高维、动态且难以精确建模的决策过程。传统基于固定规则或静态假设的调度机制在复杂环境下逐渐暴露出适应性不足的问题。

现有操作系统的调度机制普遍建立在 SMP 模型之上，假设各计算核心在性能与能耗方面等价。这种假设在经典与现代系统中均有体现：教学操作系统 xv6 的调度循环在每个 CPU 核心上独立运行，所有核心采用相同策略；Zephyr 与 Linux 内核的 CFS（Completely Fair Scheduler）均基于 per-CPU 运行队列和负载均衡机制以实现“公平”分配，默认核心性能一致。FreeRTOS 在文档中虽提及可运行于多核场景（AMP或SMP），但其标准实现与示例主要针对单核或对称多核平台，未涉及异构核性能差异或异构感知调度策略。这些系统的调度算法——如先来先服务（FCFS）、最短作业优先（SJF）、时间片轮转（RR）——在设计时均假设\textbf{任务可独立、核心可对称、环境可静态}。

然而在异构体系下，性能差异、能耗约束、任务依赖与迁移成本交织，调度问题呈现出高度耦合的动态性。
传统调度机制因此暴露出三类局限：
\begin{itemize}
    \item 缺乏动态反馈：策略在运行期无法根据系统状态自我调整；
    \item 局部最优陷阱：多核间负载均衡基于局部信息，难以形成全局最优分配；
    \item 规则僵化：调度逻辑以固定启发式为核心，无法通过经验积累改进策略。
\end{itemize}

随着异构平台普及，这些假设的崩塌直接推动了从启发式调度向自学习调度的转变。从研究演进的角度看，调度算法的变化与问题结构的演化高度一致：
\begin{table}[htbp]
\centering
\caption{调度问题在异构计算环境中的演化}
\begin{tabular}{cccc}
\hline
问题维度 & 传统阶段 & 异构阶段 & 对应算法演化 \\
\hline
时间维度 & 静态（一次性决策） & 动态（持续交互） & 启发式 $\rightarrow$ 自适应学习 \\
函数结构 & 线性、可分解 & 非线性、耦合 & 优化算法 $\rightarrow$ 强化学习 \\
资源结构 & 对称同构 & 异构非对称 & 负载均衡 $\rightarrow$ 策略自适应 \\
反馈特征 & 明确、即时 & 延迟、非确定 & 单步最优 $\rightarrow$ 序贯最优 \\
\hline
\end{tabular}
\end{table}
在传统线性结构下，算法依赖显式目标函数与可解析约束；而在动态非线性环境中，状态空间呈指数增长，决策过程本质上成为一个部分可观测、延迟反馈的序贯决策问题。这正是强化学习（Reinforcement Learning, RL）所擅长的范畴：通过“状态–动作–奖励”循环在环境交互中逐步逼近最优策略。相比之下，监督学习依赖先验标签，而调度任务中不存在可预定义的“最优决策”；无监督学习只能挖掘模式，无法实现目标驱动的策略优化。因此，强化学习在结构上天然契合异构调度问题的特征：无标签、连续反馈、长期优化、环境自适应。

\subsection{国内外研究现状}
\subsubsection{传统调度算法研究}
早期操作系统主要运行在单处理器环境中，因此调度算法设计相对简单，主要采用基于规则的经典调度策略。

在传统调度算法中，最常见的方法包括先来先服务（First Come First Served, FCFS）、最短作业优先（Shortest Job First, SJF）以及时间片轮转（Round Robin, RR）等。FCFS 按照进程到达的先后顺序进行调度，算法实现简单，但容易产生“长作业阻塞短作业”的问题。SJF 则优先调度预计运行时间较短的任务，可以在理论上最小化平均等待时间，但在实际系统中难以准确预测任务运行时间。时间片轮转算法通过为每个进程分配固定时间片，使得多个任务能够轮流获得 CPU，从而提升系统的交互性，因此被广泛应用于分时操作系统中。

随着操作系统应用场景的复杂化，研究者进一步提出了优先级调度（Priority Scheduling）以及多级反馈队列调度（Multi-Level Feedback Queue, MLFQ）等算法。优先级调度通过为不同任务设置优先级，使关键任务能够获得更多计算资源，但可能导致低优先级任务长期得不到执行。MLFQ 则通过多个队列以及动态优先级调整机制，使得系统能够在吞吐量和响应时间之间取得更好的平衡，因此成为许多操作系统调度设计的重要参考。

在现代操作系统中，调度算法进一步向公平性与可扩展性方向发展。例如 Linux 操作系统采用完全公平调度器（Completely Fair Scheduler, CFS），通过引入虚拟运行时间（virtual runtime）的概念，使各个进程能够按照相对公平的方式共享 CPU 资源。CFS 使用红黑树维护可运行进程集合，从而在保证公平性的同时维持较高的调度效率。

总体而言，传统操作系统调度算法在规则明确、实现简单的系统环境中表现良好，但这些算法通常依赖固定策略，难以根据系统运行状态进行动态调整。在多核处理器、异构计算架构以及复杂负载环境下，如何根据系统状态自适应地调整调度策略成为新的研究问题，这也推动了强化学习等智能方法在调度问题中的研究与应用。
\subsubsection{强化学习在调度中的应用}
随着计算系统规模和复杂度的不断提升，传统基于固定规则或启发式策略的调度算法在动态环境中的适应能力逐渐受到限制。近年来，强化学习（Reinforcement Learning, RL）因其能够通过与环境交互不断优化决策策略，在动态任务调度领域受到广泛关注。强化学习通过“状态—动作—奖励”的交互机制，使调度策略能够在不断试错的过程中逐渐逼近最优方案，从而实现系统资源分配与任务执行顺序的优化。

在强化学习框架下，任务调度问题通常被建模为马尔可夫决策过程（Markov Decision Process, MDP）。系统当前的任务队列状态、资源使用情况以及系统负载等信息构成环境状态，调度器根据当前状态选择相应的调度动作，例如选择某一任务进行执行或调整任务优先级。系统根据调度结果给出奖励反馈，例如减少平均响应时间或提高系统吞吐率。通过不断更新策略函数或价值函数，调度代理能够逐渐学习到适合当前系统环境的调度策略。

在具体算法方面，强化学习调度方法主要包括基于价值函数的方法、基于策略的方法以及深度强化学习方法等。其中，Q-Learning 是较早应用于调度问题的强化学习算法之一，该方法通过更新状态—动作价值函数来逐步逼近最优策略，具有实现简单、计算开销较低等优点。随着深度学习技术的发展，研究者进一步提出将深度神经网络与强化学习结合形成**深度强化学习（Deep Reinforcement Learning, DRL）**方法，例如 Deep Q-Network（DQN）。该类方法能够利用神经网络对高维状态空间进行表示，从而适用于更加复杂的调度环境，例如多任务并发执行、异构计算资源管理以及大规模计算集群调度等场景。

在操作系统任务调度层面，Sun 等人（2025）在 《Dynamic Operating System Scheduling Using Double DQN》 中首次将深度强化学习引入操作系统内核调度，构建基于 Double DQN 的调度算法，能在不同负载条件下自适应调整任务优先级与资源分配，从而提升系统吞吐率与响应速度。实验表明，该方法在轻载、中载、重载条件下均优于传统 RR 与 SJF 算法，展现出在多任务环境下的泛化能力。Chen 等人（2022）在 《Smart Scheduler: An Adaptive NVM-Aware Thread Scheduling Approach on NUMA Systems》 中提出基于扩展 LinUCB 的强化学习调度算法，用于 NUMA 架构下 NVM 访问不均衡问题。该算法结合操作系统事件与访问特征实现线程迁移自适应决策，在实际 NUMA 服务器上验证其性能提升，程序执行时间减少 59.9\%，收敛速度在 20 次动作内完成。此外， Shen 等人（2024）在 《Reinforcement Learning-Based Task Scheduling for Heterogeneous Computing in End-Edge-Cloud Environment》 中提出了基于 Q-Learning 的分层调度策略，通过结合任务特征与节点负载，实现端–边–云协同环境下的多层动态任务分配，显著降低平均等待与完成时间。

总体而言，强化学习为解决复杂动态调度问题提供了一种新的思路。相比传统调度算法，其优势在于能够通过在线学习不断优化调度策略，并适应不断变化的系统环境。然而，该方法在实际应用中仍面临一些挑战，例如状态空间维度较高、训练成本较大以及算法收敛稳定性问题等。因此，如何设计更加高效稳定的强化学习调度算法，仍然是当前研究的重要方向。

\subsection{研究内容与意义}
尽管已有研究验证了 RL 在动态调度中的潜力，但多数工作仍停留在用户态或仿真层，尚未深入到内核级实现。这种分离导致算法缺乏对真实内核约束（如中断、上下文切换、资源竞争）的验证，难以评估其在轻量操作系统中的可行性。此外，深度强化学习方法（如 DQN 、 A3C ）虽在高维状态下表现优越，但对算力和存储资源的需求较高，不适合嵌入式或教学型 OS 环境。在这一背景下，无模型强化学习（Model-Free RL）中的 Q-Learning 成为更具实践意义的选择：它不依赖环境转移模型，计算开销小，理论简单且可解释性强。其核心更新规则：
\begin{equation}
Q(s,a) \leftarrow Q(s,a) + \alpha \left[r + \gamma \max_{a'} Q(s',a') - Q(s,a)\right]
\end{equation}
通过与系统环境的持续交互，在有限资源下即可实现策略收敛。在操作系统调度问题中，
状态 $s$ 可定义为系统负载、任务队列与资源占用；动作 $a$ 对应于任务分配或优先级调整；
奖励 $r$ 反映吞吐率、响应时间或能耗的综合指标。这一映射使 Q-learning 能够自然嵌入
内核调度循环，实现在线自学习与策略优化。
\begin{table}[htbp]
\centering
\caption{强化学习调度相关研究对比}
\begin{tabular}{ccccc}
\hline
研究 & 算法类型 & 应用层级 & 优点 & 局限性 \\
\hline
Shyalika et al. (2020) & Q-Learning, DQN 综述 & 仿真层 & 概念框架完整 & 缺乏系统级实验 \\
Sun et al. (2025) & Double DQN & 用户态调度 & 性能优异 & 未嵌入内核 \\
Chen et al. (2022) & LinUCB & NUMA 系统 & 收敛快 & 资源模型复杂，非轻量 \\
Shen et al. (2024) & Q-Learning & 云-边-端调度 & 层次协同 & 非实时，非操作系统内核 \\
本文工作 & Q-Learning & 轻量级 OS 内核 & 低成本、自学习 & --- \\
\hline
\end{tabular}
\end{table}
基于此，本研究拟将 Q-learning 算法引入 xv6 轻量级操作系统的进程调度模块，
探索强化学习在系统级调度中的应用。通过对比不同负载下传统算法与
Q-learning 调度的性能，分析其优化效果。